\cleardoublepage
\chapter{Wstęp}
\label{cha:Wstep}
% \addcontentsline{toc}{chapter}{Wstęp}

%--------------------------------------
% początek: wskazówka
%--------------------------------------
%„Wstęp powinien przedstawiać ogólne informacje na temat, którego dotyczy praca, historię i~zakres zastosowań”~\cite{bib:ZasadyPisania}. We wstępie nie powinno być rysunków, tabel i~wzorów.
%--------------------------------------
% koniec: wskazówka
%--------------------------------------

W ostatnim czasie można zaobserować intensywny rozwój architektury procesorów wielordzeniowych. Producenci sprzętu zamiast dalszego 
zwiększania częstotliwości taktowania pojedynczego rdzenia decydują się na zwiększanie liczby rdzeni fizycznych oraz~logicznych. 
Taka zmiana podejścia projektowania procesorów stawia przed twórcami oprogramowania nowe wyzwania. Klasyczne algorytmy sekwencyjne 
nie wykorzystują w~pełni dostępnej mocy obliczeniowej współczesnych jednostek centralnych.

Szczególnie ważnym obszarem, w~którym równoległość może przynieść wymierne korzyści, jest sortowanie danych. Sortowanie stanowi 
jeden z~fundamentalnych problemów informatyki i~znajduje szerokie zastosowanie m.in. w bazach danych, systemach plików, 
przetwarzaniu dużych zbiorów danych oraz uczeniu maszynowym. Wraz ze wzrostem objętości przetwarzanych informacji klasyczne algorytmy 
sortowania sekwencyjnego mogą coraz częściej stać się wąskim gardłem wydajnościowym.

Implementacja efektywnych algorytmów sortowania równoległego wiąże się z~szeregiem wyzwań, takich jak odpowiedni podział danych,
synchronizacja jednostek wykonawczych, zarządzanie dostępem do pamięci oraz ograniczanie narzutu wynikającego z~realizacji obliczeń równoległych.
W~praktyce wiele istniejących rozwiązań albo pozostaje na poziomie teoretycznym, albo nie uwzględnia realnych ograniczeń sprzętowych 
i~systemowych, takich jak koszt przełączania kontekstu, konflikty w~dostępie do pamięci czy nierównomierne obciążenie rdzeni.

W związku z~tym praktyczne zbadanie wydajności oraz efektywności wybranych algorytmów sortowania równoległego stanowi 
istotne i~aktualne zagadnienie. W niniejszej pracy skoncentrowano się na implementacji i~porównaniu czterech popularnych 
algorytmów: Parallel~Quicksort, Parallel~MergeSort, Parallel~BucketSort oraz Parallel~Samplesort. 
Algorytmy zostały zrealizowane w~języku Python z~wykorzystaniem biblioteki \texttt{multiprocessing}, a~następnie poddane testom 
pod kątem czasu wykonania, wykorzystania pamięci RAM, obciążenia procesora, skalowalności oraz wpływu charakterystyki danych wejściowych.

\section{Cel pracy}
\label{sec:CelPracy}
% \addcontentsline{toc}{section}{Cel pracy}

Głównym celem niniejszej pracy magisterskiej jest przeprowadzenie szczegółowej analizy wydajności i~efektywności wybranych algorytmów 
sortowania równoległego na współczesnych wielordzeniowych procesorach. Praca skupia się na praktycznej implementacji oraz 
porównaniu czterech algorytmów: Parallel Quicksort, Parallel MergeSort, Parallel BucketSort oraz Parallel Samplesort.

Kluczowym celem jest zbadanie, w~jaki sposób wybrane algorytmy zachowują się w~realnych warunkach obliczeniowych, a~nie 
jedynie w~warunkach idealnych zakładanych przez modele teoretyczne. W~tym celu algorytmy zostały zaimplementowane w~języku 
Python z~wykorzystaniem biblioteki \texttt{multiprocessing}, która umożliwia tworzenie procesów oraz współdzielenie pamięci. 
Każdy z~algorytmów został uruchomiony zarówno w~wariancie sekwencyjnym stanowiącym punkt odniesienia, jak i~w wariancie równoległym, 
przy różnej liczbie wykorzystywanych rdzeni procesora.

Badanie obejmuje kilka kluczowych aspektów:

\begin{itemize}
    \item pomiar czasu sortowania w~zależności od rozmiaru zbioru danych oraz liczby użytych rdzeni,
    \item analizę zużycia pamięci operacyjnej (RAM) oraz obciążenia procesora (CPU) podczas wykonywania kodu,
    \item ocenę skalowalności, czyli zdolności algorytmów do efektywnego wykorzystania rosnącej liczby rdzeni,
    \item porównanie zachowania się algorytmów przy różnych charakterystykach danych wejściowych tj. dane losowe, częściowo posortowane 
    oraz z duplikatami.
\end{itemize}

Dodatkowym celem pracy jest wyznaczenie i~analiza klasycznych metryk równoległości, \textit{speedup} (przyspieszenie) oraz 
\textit{efficiency} (efektywność), a~także identyfikacja głównych czynników ograniczających wydajność poszczególnych algorytmów 
(m.in. narzut związany z~tworzeniem procesów, synchronizacją, komunikacją oraz zarządzaniem pamięcią współdzieloną).

W rezultacie pracy oczekuje się uzyskania praktycznych wniosków dotyczących tego, które z~badanych algorytmów sortowania równoległego 
najlepiej sprawdzają się w~środowisku wielordzeniowym przy określonych rozmiarach i~typach danych, a~także wskazania ich mocnych 
i~słabych stron w~kontekście rzeczywistego wykorzystania zasobów obliczeniowych.

\section{Zakres pracy}
\label{sec:ZakresPracy}
% \addcontentsline{toc}{section}{Zakres pracy}

Zakres pracy obejmuje:

\begin{itemize}
    \item zebranie wiadomości z~zakresu sortowania równoległego, architektury procesorów wielordzeniowych, metryk wydajności 
    i~efektywności oraz skalowalności algorytmów równoległych,

    \item porównanie wybranych algorytmów sortowania równoległego (Parallel Quicksort, Parallel MergeSort, Parallel BucketSort oraz 
    Parallel Samplesort) pod kątem wydajności, skalowalności i~zużycia zasobów obliczeniowych,

    \item opracowanie założeń projektowych dotyczących implementacji algorytmów w~środowisku wielordzeniowym z~wykorzystaniem 
    biblioteki \texttt{multiprocessing} oraz pamięci współdzielonej,

    \item implementację wariantów sekwencyjnych i~równoległych wybranych algorytmów sortowania w~języku Python,

    \item implementację generatora danych testowych o~różnych charakterystykach (dane losowe, częściowo posortowane, oraz dane z duplikatami),

    \item przetestowanie zaimplementowanych algorytmów na wielordzeniowych procesorach, obejmujące pomiar czasu wykonania, zużycia pamięci RAM, 
    obciążenia CPU, wyznaczenie metryk \textit{speedup} i~\textit{efficiency} oraz analizę wpływu charakterystyki danych wejściowych.

    \item prezentację wyników badań za pomocą tabel oraz wykresów.
\end{itemize}

------------------- Do uzupełnienia na końcu pracy -------------------

W rozdziale~\ref{cha:RozdzialWprowadzajacy} przedstawiono podstawy teoretyczne sortowania równoległego, 
charakterystykę procesorów wielordzeniowych oraz metryki wykorzystywane do oceny wydajności i~skalowalności.  
W rozdziale~\ref{cha:WybraneAlgorytmy} opisano wybrane algorytmy sortowania równoległego.  
W rozdziale~\ref{cha:SrodowiskoImplementacja} scharakteryzowano środowisko badawcze, zastosowane technologie oraz 
sposób implementacji algorytmów.  
W rozdziale~\ref{cha:MetodykaBadan} przedstawiono metodykę badań, scenariusze testowe oraz narzędzia pomiarowe.  
W rozdziale~\ref{cha:AnalizaBadan} zaprezentowano i~przeanalizowano wyniki przeprowadzonych eksperymentów.  
Pracę zamyka podsumowanie zawierające najważniejsze wnioski oraz możliwe kierunki dalszych badań.

-----------------------------------------------------------------------